\section{Finding 2: What ORB Costs in Storage and Time}
\label{sec:storage}

The paper's design stores every image's hashes inside its mint transaction, so that detection
is ``fully self-contained within the blockchain''. That works for the four hashes. It does not
work for ORB, and we report this as a finding.

\begin{table}[ht]
\centering
\begin{tabular}{@{}lrr@{}}
\toprule
Representation & Bytes per image & Transactions (400\,B each) \\
\midrule
aHash + pHash + hsvHash          & 24 (fixed)    & 0.06 \\
sHash                            & 32 (median)   & 0.08 \\
\textbf{All four paper hashes}   & \textbf{$\approx$56} & \textbf{0.14} \\
\midrule
ORB, default (about 455 landmarks $\times$ 32\,B) & 14,560 & 36.4 \\
\bottomrule
\end{tabular}
\caption{Measured storage per image, on 60 real originals.}
\label{tab:storage}
\end{table}

At OpenCV's default setting, ORB keeps about 455 descriptors of 32 bytes each, 14,560 bytes per
image, the size of about 36 whole transactions. The default cap of 500 descriptors was never a
deliberate choice, and it is not even reached. So we measured the whole accuracy-versus-size
curve, giving every descriptor budget its own best threshold so that only the shape of the
curve is compared (Figure~\ref{fig:storage}). This protocol checks out: at the default budget
it gives F1 90.8, within 0.2 of the train-tuned 90.6.

\begin{figure}[ht]
\centering
\includegraphics[width=0.72\linewidth]{\fig{fig_storage.pdf}}
\caption{ORB accuracy on the geometric subset against bytes stored per image. The shaded band
is one mint transaction. No point on the ORB curve lies inside it.}
\label{fig:storage}
\end{figure}

\begin{table}[ht]
\centering
\begin{tabular}{@{}rrrrrr@{}}
\toprule
Landmarks & Bytes & Transactions & Precision & Recall & F1 \\
\midrule
455 (default) & 14,560 & 36.4 & 90.8\% & 90.8\% & 90.8\% \\
242           &  7,760 & 19.4 & 94.2\% & 86.1\% & 89.9\% \\
122 (cap 128) &  3,888 &  9.7 & 96.1\% & 80.2\% & 87.5\% \\
64            &  2,048 &  5.1 & 94.8\% & 74.7\% & 83.5\% \\
32            &  1,024 &  2.6 & 90.8\% & 65.5\% & 76.1\% \\
\midrule
sHash (oracle) & 32 & 0.08 & 66.5\% & 89.8\% & 76.4\% \\
\bottomrule
\end{tabular}
\caption{Selected points of the size curve (geometric subset, best threshold per budget).}
\label{tab:budget}
\end{table}

\begin{itemize}
  \item \textbf{128 landmarks is the practical sweet spot:} 3,888 bytes (about 10 transactions)
        for F1 87.5, still far above sHash. As the budget shrinks, precision rises and recall
        falls; ORB becomes a stricter matcher.
  \item \textbf{32 landmarks is the floor:} 1,024 bytes, still 2.6 transactions, and F1 76.1,
        which only ties sHash (76.4) at 32 times its size. Below this, the best threshold
        collapses to ``flag everything''.
\end{itemize}

\paragraph{Conclusion.} There is no size at which ORB both fits in a transaction and beats
sHash. The limit is architectural, not a matter of tuning: a signal made of a set of local
features cannot travel in transaction metadata. The descriptor type also decides the index:
one fixed-width string fits a BK-tree, while a set of binary descriptors needs LSH. The
signal and its storage cannot be chosen separately. Section~\ref{sec:future} discusses ways
around the limit.

\subsection{Time}
\label{sec:time}

Storage asks what ORB costs to keep. We also measured what it costs to run, on the two tasks
both signals perform: \emph{fingerprinting a new image} (ORB finds and describes its
landmarks; sHash segments the image into its main objects and hashes each one) and
\emph{comparing one pair} (ORB matches the two landmark sets and fits a geometric transform
with RANSAC, once for each of the four mirror variants; sHash computes the paper's
mean-of-minimums distance). Search and index time were not measured for either signal.

\paragraph{Two comparisons, and why.} Section~\ref{sec:goals} explains that a C11-versus-Python
speed claim is not fair: it mixes language effects with algorithmic ones. We therefore report
two rows and say what each answers. The \emph{as-used} row is the tools as actually built:
\texttt{imagehash}, the Python library behind the paper's numbers, against OpenCV's ORB (C++).
It mixes languages and is not an algorithm comparison. The \emph{like-for-like} row is the
algorithm comparison: our bit-exact C11 port of sHash (Section~\ref{sec:reimpl}) against
OpenCV's ORB, both compiled code. Even there the remaining bias runs against sHash, since
OpenCV is vectorised and our port is a plain \texttt{-O2} build.

\paragraph{Protocol.} 1,168 images and 800 pairs from the test split (100 pairs of every edit
type, non-copies included); one warm-up and five timed runs per item, keeping the median; a
single thread on an Apple M5. Decoding the image file is excluded on both sides: both signals
need the pixels, and each side decodes with a different library (Pillow takes about 1.1\,ms
per image). The timed code is the evaluated code: every timed sHash distance equals the one
used for accuracy, and the C port, split into decode plus hash for this purpose, still matches
\texttt{imagehash} bit for bit on all 12,000 test images. Two full runs agree to within 3\%.

\begin{table}[ht]
\centering\small
\begin{tabular}{@{}lrrrr@{}}
\toprule
Median per item & sHash, as used & sHash, like for like & ORB & ORB \\
                & (\texttt{imagehash}) & (C11 port) & (455 landmarks) & (128) \\
\midrule
Fingerprint a new image & 67.9\,ms & \textbf{1.46\,ms} & \textbf{0.94\,ms} & 0.72\,ms \\
\quad ORB query side ($\times$4 mirrors) & & & 3.97\,ms & 3.12\,ms \\
\textbf{Compare one pair} & 11.3\,\textmu s & \textbf{0.1\,\textmu s} & \textbf{34.5\,ms} & 24.3\,ms \\
\bottomrule
\end{tabular}
\caption{Time per item, test split. The as-used column mixes Python with C++; the
like-for-like columns are the algorithm comparison. ORB at 128 landmarks is the size curve's
sweet spot (Table~\ref{tab:budget}).}
\label{tab:time}
\end{table}

\begin{figure}[ht]
\centering
\includegraphics[width=0.72\linewidth]{\fig{fig_time.pdf}}
\caption{The like-for-like rows of Table~\ref{tab:time}, on ordinary millisecond axes. The two
signals take about as long to fingerprint an image; sHash's pair comparison (0.1\,\textmu s) is
too short to draw next to ORB's.}
\label{fig:time}
\end{figure}

\begin{itemize}
  \item \textbf{Making a fingerprint is not where the two differ.} Like for like, ORB takes
        0.94\,ms and sHash 1.46\,ms. The as-used 68\,ms is mostly language: about 70\% of it
        is the pure-Python flood fill inside \texttt{imagehash}.
  \item \textbf{Comparing a pair is where they differ:} about 34\,ms for ORB against about
        0.1\,\textmu s for sHash (11\,\textmu s in Python). Two fixed-width hashes are compared
        with a handful of bit counts; two landmark sets need a matching step and a RANSAC fit,
        four times over. At this per-pair cost, ORB needs an index such as LSH to scale.
  \item \textbf{A smaller ORB saves space, not much time.} Matching the descriptors is only
        about 5.9 of the 34.5\,ms; RANSAC is the rest. At 128 landmarks a pair still takes
        24.3\,ms, although the bytes fall by 74\%.
  \item \textbf{At full resolution} (a sensitivity check on 50 original images per collection),
        fingerprinting a 2000\,px Azuki image takes 16.3\,ms for ORB, about 14.7\,ms of it the
        shrink to 256\,px, against 6.8\,ms for the C port of sHash. Pair times do not depend on
        resolution.
\end{itemize}

The paper reports its own search times (its Table~VII: from 120\,ms over 100 stored images to
about 11\,s over a million, for its Python module with four BK-trees). Those time a whole search
on different hardware, while ours time single operations, so we cite them as context only.
